% Include after amsmath, amssymb, amsthm, booktabs, and xurl.
% Ordinary written arithmetic proof; not a new Lean development.
\section{A relative iPRA proof for the reassociated operand reader}
\label{sec:orh-relative}

Write $M_0$ for the revised packet's literal PA194AX11H machine and
$M_1$ for PA194ORH.  Both have 194 working controls, two symbols,
left/right moves, a separate halt control, and a blank initial tape.
The following argument proves their blank-tape halting equivalence in
intuitionistic primitive recursive arithmetic.  Composing it with the
revised AX11H written calibration yields the same consistency endpoint
for $M_1$.  This section supplies an ordinary written proof with exact
finite certificates; it does not assert a new end-to-end Lean theorem.

The dense table identities, in the order $M_0,M_1$, are
\begin{center}\footnotesize
\begin{tabular}{l}
\path{c5dfe1026d077d3e062eca11edb789597bfbbced4f6040f8828ae5d467b1167a}\\
\path{c16bef7a1587e31adaae747d5bad26721f30bb5c6c33dd4b798be95f5141f1d4}.
\end{tabular}
\end{center}
The native tables start at control 74, use 84 as the source boundary,
and halt at 254.  Their supplied injective label maps identify these
with dense start 0 and halt 194.  The label maps are conjugacies of
every literal transition, so quantifier-free induction on time transfers
the statements below to the dense tables.

\subsection{The changed reader and its second entry}

In the source tree put
\[
 A=\texttt{node32},\quad B=\texttt{node53},\quad
 C=\texttt{node55},\quad D=\texttt{node65}.
\]
The sole named-tree reassociation is
\[
 \texttt{node75}=((A,(B,C)),D)
       \quad\longmapsto\quad (A,(B,(C,D))).
 \tag{OR1}
\]
Both trees have the same left-to-right leaf order.  Their exact
operand-reader codes and exit controls are:
\begin{center}
\begin{tabular}{cccc}
\toprule
Operand & Old word & New word & Native exit\\
\midrule
$A$ & 00 & 0   & 104\\
$B$ & 010& 10  & 107\\
$C$ & 011& 110 & 98\\
$D$ & 1  & 111 & 90\\
\bottomrule
\end{tabular}
\end{center}
These are complete prefix codes.  Starting at control 213, each row
preserves its scanned bit and moves right.  The new rows are
\[
\begin{array}{c|cc}
q&\text{on }0&\text{on }1\\\hline
213&0R104&1R185\\
185&0R107&1R333\\
333&0R98&1R90.
\end{array}
\tag{OR2}
\]
Control 333 also supplies the shared $(C,D)$ reader.  The old reader
used controls 213, 92, 185, with respective rows
$(0R92,1R90)$, $(0R104,1R185)$, and $(0R107,1R98)$.

There is an additional entry into the old control 213 from the history
selector.  It must be preserved separately.  The new rows are
\[
 \delta_1(69,0)=0L92,\qquad
 \delta_1(92,0)=0R111,\qquad
 \delta_1(92,1)=1R90.
 \tag{OR3}
\]
The unchanged row at $(111,0)$ is $0R104$.
Suppose control 69 scans a zero at coordinate $h$, and let $b$ be
the bit at $h-1$.  For $b=0$ the old and new control paths are
\[
 69\longrightarrow213\longrightarrow92\longrightarrow104,
 \qquad
 69\longrightarrow92\longrightarrow111\longrightarrow104.
\]
Both take three steps, preserve the two zero cells, and finish with
the head at $h+1$.  For $b=1$ the respective paths are
$69\to213\to90$ and $69\to92\to90$; both take two steps and
finish at $h$, preserving both cells.  Neither computation reads
outside these two cells, so this is a uniform identity for arbitrary
surrounding tapes.  The row $(69,1)$ is unchanged.  Thus (OR3)
preserves the second entry while (OR2) implements the new operand code.

The exact resolved primitive instruction arrays are byte-identical:
3,671 sites and 4,778 outcomes, with both outcomes of every decrement
included.  Both arithmetic and interval-based source resolvers produce
that same array.  There are 134 changed source addresses.  This finite
graph equality, rather than an assumption that arbitrary tree
reassociations preserve jump depths, identifies the common source
program used below.  Both initial site-25 addresses are 1000000.

\subsection{One source run, two complete tape encodings}

Use a source state
\[
 z=(i,R_0,\ldots,R_8,n,Q),\qquad Q=[s]X\ne[],
\]
or a distinguished absorbing halted state $H$.  The site $i$ ranges
over the common finite instruction array; $n$ is the nonnegative
header slack, and $Q$ is a nonempty finite list of natural numbers.
Define $F$ by the common instruction array, with the following
bookkeeping.  An increment appends the value 1 to $Q$.  A positive
decrement increases $n$ by one.  Other live instructions leave $n,Q$
unchanged.  A halt instruction yields $H$, and $F(H)=H$.
All instructions and list operations have primitive recursive codes.

For $j=0,1$, let $p_j(i)$ be the corresponding finite source-address
lookup.  Put
\[
 \pi=(6,2,0,4,3,7,8,5,1),\qquad V_{\pi(a)}=R_a,
 \qquad B(v)=0\,1^{v+1}.
\]
The complete boundary encoding $E_j(z)$ has control 84, head at
$-2$, outward left stream $1$ followed by infinite zeros, and right
stream
\[
 p_j(i)\;0^{48+n-|p_j(i)|}\;1\;
 B(V_0)\cdots B(V_8)B(s)B(X)\;0^\infty.
 \tag{OR4}
\]
All addresses have length at most 37.  Consequently (OR4) is defined
for every valid source state, including zero registers and empty
tail $X$.  Set $E_j(H)=H$.

Finite binary tape halves can be encoded by natural numbers $L,R$,
with the scanned bit the low bit of $R$.  Writing $w$ and moving
right sends them to $(2L+w,\lfloor R/2\rfloor)$; moving left sends
them to
\[
 (\lfloor L/2\rfloor,
  (L\bmod2)+2w+4\lfloor R/2\rfloor).
\]
After reaching the halt control, replace the tape configuration by
$H$ and keep it fixed.  This arithmetic convention preserves the
literal machine's halt predicate, by induction on time.  Write
$\operatorname{Run}_j(t,c)$ for the primitive recursive iterate.
The semigroup equation
\[
 \operatorname{Run}_j(u+v,c)
  =\operatorname{Run}_j(v,\operatorname{Run}_j(u,c))
 \tag{OR5}
\]
and absorption of $H$ have quantifier-free induction proofs.

\subsection{Uniform local contracts and their exact durations}

Each outcome certificate records a 13-entry nonnegative coefficient
vector.  Its valuation is
\[
 (1,n,x_0,\ldots,x_8,s,W(X)),\qquad
 W(X)=\sum_{v\in X}(v+2).
 \tag{OR6}
\]
For a positive decrement of $R_a$, set $x_a=R_a-1$ and use
$R_a=x_a+1$; the other slots equal their logical registers.
All other outcome valuations use their stated zero test or register
values.  Crucially, the final slot is the width of the very same
list $X$ in (OR4), not an independent free parameter.
Select the exact duration $D_j(z)$ by the source site and decidable
decrement branch, and the corresponding scalar product (OR6).
Set $D_j(H)=0$.

The full independent literal replays accept all 4,778 outcomes for
each table.  For $M_1$ they contain 1,102,204 symbolic events, with
at most 403 events in any one outcome.  They use literal transitions,
unary self-loops, normalization by finite-word identities, and four
literal block traversals.  No arithmetic-box oracle is used.

Here is why those certificates give uniform arithmetic theorems.
A unary self-loop crossing a run of length $a$ is verified by
quantifier-free induction on the processed length.  Each block
traversal takes exactly $v+2$ steps on $B(v)$ or its specified
intermediate block; induction on a fixed list's prefix proves a
whole-list duration $W(X)$ with both exterior contexts untouched.
The normalization identities
\[
 C(X)1=1B(X),\qquad
 1\operatorname{rev}(C(X))=\operatorname{rev}(B(X))1,
 \quad C(v)=10\,1^v,
\]
have the same prefix-induction proof.  Literal events are direct
substitutions in the tape-step equations.  Induction on a successful
event trace, using (OR5), therefore proves its exact complete-tape
endpoint at its recorded affine duration.  The invariant is a
bounded Boolean check and an equality of PR configuration functions.
It contains no unbounded quantifier over environments.

The four block entries and rows are unchanged, and were rederived
from both tables.  Every other changed row is checked as a literal
transition.  Hence the event-calculus soundness proof in the revised
AX11H packet applies to the new successful traces as well.
All 4,777 nonhalting endpoints have zero head displacement and
zero address overflow.  Independently checking their computed
addresses against the common source graph gives
\[
 p_{\rm computed}0^{48-|p_{\rm computed}|}
   =p_j(i')0^{48-|p_j(i')|}.
 \tag{OR7}
\]
Appending the same extra slack zeros proves closure of (OR4).
The single halt outcome has exact duration 3.  These finite checks
and the event-calculus induction establish in iPRA, for every valid
source state,
\[
 \operatorname{Run}_j(D_j(z),E_j(z))=E_j(F(z)),\qquad
 z\ne H\ \Longrightarrow\ D_j(z)\ge3.
 \tag{OR8}
\]
The source validity predicate is preserved.  At $H$, (OR8) is the
identity $\operatorname{Run}_j(0,H)=H$; no positive duration is
asserted for the halted state.

\subsection{Explicit primitive recursive clocks and a polynomial bound}

The common initial source state is
\[
 z_*=(25,2,0,\ldots,0,0,[1]^{26}).
\]
Independent literal startup runs establish the complete tape
equalities
\[
 \operatorname{Run}_j(T_j,\mathrm{blank})=E_j(z_*),
 \qquad T_0=57922,\quad T_1=57942.
 \tag{OR9}
\]
The head remains inside the checked finite array throughout each
run, all exterior cells are untouched zeros, and the entire final
array equals (OR4).  Thus (OR9) is a closed finite arithmetic fact,
not a comparison of a few selected tape cells.

Define by primitive recursion
\[
 z_k=F^k(z_*),\qquad
 \tau_j(0)=0,\qquad
 \tau_j(k+1)=\tau_j(k)+D_j(z_k).
 \tag{OR10}
\]
Quantifier-free induction, with source validity included in the
invariant, proves
\[
 \operatorname{Run}_j(T_j+\tau_j(k),\mathrm{blank})=E_j(z_k).
 \tag{OR11}
\]
This equation remains exact after halt: both source and target are
absorbing and subsequent durations are zero.

For both tables the componentwise maximum of the affine duration
vectors is exactly
\[
 (534,8,4,4,4,4,4,4,2,4,4,2,2).
 \tag{OR12}
\]
Consequently each live source state's duration satisfies
\[
 3\le D_j(z)\le534+8n+4\sum_{a=0}^8R_a+2W(Q).
 \tag{OR13}
\]
For the two spare slots, $s+W(X)\le W(Q)=s+2+W(X)$;
for a positive decrement its variable slot is at most $R_a$.
These justify (OR13) directly from the recorded coefficients.

Along the particular run from $z_*$, another quantifier-free
induction gives, while live,
\[
 n_k\le k,\qquad \sum_aR_{a,k}\le2+k,
 \qquad W(Q_k)\le78+3k.
 \tag{OR14}
\]
Indeed, one source instruction increments at most one register;
only a positive decrement increases slack, by one; and an increment
appends one block $B(1)$ of width three.  Initially the register
sum is two and the 26 spare blocks have total width 78.
Assigning zero to these auxiliary selectors at $H$ extends the
bounds across canonical halt.  Hence
\[
 0\le D_j(z_k)\le698+18k,\qquad
 \tau_j(k)\le698k+18\sum_{r<k}r=9k^2+689k.
 \tag{OR15}
\]
All of these are PR arithmetic identities or inequalities with
quantifier-free induction invariants.  Define the explicit
polynomial clocks
\[
 P_0(k)=57922+9k^2+689k,\qquad
 P_1(k)=57942+9k^2+689k.
 \tag{OR16}
\]
The exact clocks are $T_j+\tau_j(k)$; $P_j$ are sufficient
uniform upper bounds at which absorption preserves any halt
already reached.

\subsection{Halting reflection without positive stopped durations}

From (OR8), source absorption, and quantifier-free induction,
\[
 z_k\ne H\quad\Longrightarrow\quad\tau_j(k)\ge3k.
 \tag{OR17}
\]
For the successor case, $z_{k+1}\ne H$ implies $z_k\ne H$,
so the next duration adds at least three.  No global cofinality
of $\tau_j$ after halt is claimed or needed.

Suppose $M_j$ has halted by physical time $t$.  Put $k=t+1$.
If $z_k\ne H$, then (OR17) implies
$T_j+\tau_j(k)>t$.  Target absorption makes the left side of
(OR11) halted, but its right side is a live frame at control 84.
This is a contradiction.  Equality of the source state with its
halt code is decidable, so intuitionistic logic gives $z_{t+1}=H$.
Conversely, $z_k=H$ gives physical halt at $T_j+\tau_j(k)$
by (OR11), and therefore by $P_j(k)$ by (OR15) and absorption.

Write $h_j(t)$ for the atomic PR assertion that $M_j$ is halted
by time $t$ from blank tape.  We have proved the pointwise witness
transports
\begin{align*}
 \mathrm{iPRA}&\vdash h_0(t)\ \longrightarrow\ h_1(P_1(t+1)),\\
 \mathrm{iPRA}&\vdash h_1(t)\ \longrightarrow\ h_0(P_0(t+1)).
 \tag{OR18}
\end{align*}
These are explicit PR functions in both directions.  Their proof
uses neither an unbounded least-time search nor Markov's principle.

\paragraph{Relative theorem.}
For the exact two literal tables identified above,
\[
 \mathrm{iPRA}\ \vdash\
 \operatorname{Halt}_{\rm blank}(M_0)
   \longleftrightarrow\operatorname{Halt}_{\rm blank}(M_1),
 \tag{OR19}
\]
and also
\[
 \mathrm{iPRA}\ \vdash\
 \neg\operatorname{Halt}_{\rm blank}(M_0)
   \longleftrightarrow\neg\operatorname{Halt}_{\rm blank}(M_1).
 \tag{OR20}
\]
\begin{proof}
Existential introduction and elimination applied to (OR18) give
(OR19).  To derive, for example, nonhalting of $M_1$ from nonhalting
of $M_0$, assume a halting witness for $M_1$, transport it by the
second implication of (OR18), and contradict nonhalting of $M_0$.
The opposite implication is the same use of the first transport.
This is intuitionistic reasoning.  The equivalent notation
$\forall t\,\neg h_j(t)$ uses the intuitionistically valid
equivalence $\neg\exists t\,h_j(t)\leftrightarrow\forall t\,\neg h_j(t)$.
\end{proof}

\subsection{The absolute consistency corollary and its provenance}

The exact revised AX11H calibration is retained as
\path{reference/calibration_AX11H.zip}, of SHA256
\begin{center}\footnotesize
\path{93ee082ef035b2d32e71d6af5239b25b4bea1773f03bea4d2c39810a7fe15075}.
\end{center}
Its written result for $M_0$ is
\[
 \mathrm{iPRA}\vdash
 \operatorname{Con}(\mathrm{PA})
   \longleftrightarrow\neg\operatorname{Halt}_{\rm blank}(M_0).
 \tag{OR21}
\]
Here $\operatorname{Con}(\mathrm{PA})$ is the particular mapped native
Foundation PA consistency sentence used in that calibration, whose
current Foundation endpoint is bound explicitly in its proof.
The packet supplies the source-command theorem, direct numerical
proof-list compilation, the current Foundation-to-ordinary-PAeq
adapter, and the model-uniform weak-arithmetic transfer.  The review
\path{comparison/calibration_review.txt} describes the inspected
proof strengths and the unfinished Lean composition.

\paragraph{Absolute corollary.}
Using that written AX11H theorem, the present construction proves
\[
 \boxed{\quad\mathrm{iPRA}\vdash
 \neg\operatorname{Halt}_{\rm blank}(\mathrm{PA194ORH})
   \longleftrightarrow\operatorname{Con}(\mathrm{PA}).\quad}
 \tag{OR22}
\]
\begin{proof}
Compose (OR20) with (OR21), in the indicated direction of each
implication.  The composition is intuitionistic and introduces no
new principle of arithmetic or consistency assumption.
\end{proof}

The finite inspector \path{relative_ORH_check.py} verifies the
common source array, the two history-selector cases, both prefix
readers, every padded source endpoint, the complete contract index,
the affine maxima (OR12), and the constants in (OR16).  Its receipt
is \path{relative_ORH_inspection.json}.  The literal replay supplies
the successful traces and complete startup binding; the arithmetic
induction arguments above give their uniform meaning.  These are
separate, complementary parts of the proof.
